\section{When does an FFM beat a GBDT?}
\label{sec:ffmvsgbdt}

Before adding relational structure we establish that the FFM advantage is
\emph{dataset-dependent}, which motivates the relational question. On IBM TabFormer (synthetic
card transactions with per-account injected fraud), a converged FFM (frozen probe) reaches
PR-AUC \num{0.807} versus \num{0.369} for LightGBM on matched causal features --- a $2\times$
improvement, because TabFormer fraud is a \emph{burst process} (a compromised card fires several
fraudulent transactions within seconds) that a sequence encoder captures directly. On IEEE-CIS
(real card-not-present fraud), restricting both models to the same interpretable fields, the FFM
\emph{loses}: PR-AUC \num{0.106} vs.\ \num{0.157}, because that fraud is far more a function of the
\emph{current} transaction's attributes than of account history --- a per-row pattern trees model
directly. The FFM helps when the signal is sequential and aligned with its objective, and
underperforms trees when it is static. This sharpens the question: can we give the FFM the
\emph{cross-entity} signal a per-account model cannot see at all?
